%%-------------------------------
%% Codificação e fontes
%%-------------------------------
\usepackage[utf8]{inputenc}     
\usepackage[T1]{fontenc}        
\usepackage{mathptmx}           
\selectlanguage{brazil}         
\frenchspacing                   

%%-------------------------------
%% Matemática
%%-------------------------------
\usepackage{amsmath, amssymb, amsfonts, amsthm} 

%%-------------------------------
%% Layout e formatação
%%-------------------------------
\usepackage[bottom=2cm,top=3cm,left=3cm,right=2cm]{geometry} 
\usepackage{indentfirst}         
\setlength{\parindent}{1.25cm}  
\OnehalfSpacing                  

%%-------------------------------
%% Figuras e PDFs
%%-------------------------------
\usepackage[table,xcdraw]{xcolor}
\usepackage{graphicx,subfig,float,lscape,pdfpages} 
\graphicspath{{./figuras/}}      
\newcommand{\figfonte}[1]{\caption*{Fonte: #1}} 

%%-------------------------------
%% Código-fonte
%%-------------------------------
\definecolor{codebg}{rgb}{0.95,0.95,0.95} 
\usepackage{listings}           
\lstset{
	backgroundcolor=\color{codebg},
	basicstyle=\ttfamily\footnotesize,
	breaklines=true,
	keywordstyle=\color{blue},
	stringstyle=\color{black},
	commentstyle=\color{orange},
	language=Python,
	showspaces=false,
	showstringspaces=false,
	columns=fullflexible,
	frame=single
}

%%-------------------------------
%% Listas, colunas e Gantt
%%-------------------------------
\usepackage{multicol,enumitem,pgfgantt} 

%%-------------------------------
%% Glossário
%%-------------------------------
\usepackage[nogroupskip,nonumberlist]{glossaries} 
\makeglossaries
% \newglossaryentry{stakeholder}{
%      name=Stakeholder,
%      description = {qualquer indivíduo ou organização que, de alguma modo, é impactado pelas ações de uma determinada empresa. Em uma tradução livre para o português, o termo significa parte interessada.}
% }
  

%%-------------------------------
%% Citações ABNT
%%-------------------------------
\usepackage[abnt-emphasize=bf,abnt-and-type=e,alf]{abntex2cite} 

%%-------------------------------
%% Apêndices, diagramas e algoritmos
%%-------------------------------
\usepackage{appendix,tikz}      
\usetikzlibrary{petri, positioning} 
\usepackage[ruled,vlined]{algorithm2e} 

%%-------------------------------
%% Seções e subseções
%%-------------------------------
\setsecheadstyle{\bfseries \normalsize \uppercase}  
\setsubsecheadstyle{\normalsize \uppercase}        
\setsubsubsecheadstyle{\bfseries \normalsize}      

%%-------------------------------
%% Hiperlinks e ORCID
%%-------------------------------
\usepackage{hyperref}           
\hypersetup{colorlinks=true, linkcolor=blue, citecolor=blue!85} 
\usepackage{orcidlink}          

%%-------------------------------
%% Cabeçalho com seção e pagestyle
%%-------------------------------
\usepackage{fancyhdr}

%%-------------------------------
%% Substituir chapter por section
%%-------------------------------
\let\oldchapter\chapter
\renewcommand{\chapter}[1]{\section{#1}} 

%%-------------------------------
%% Anexos: numerar seções como letras e subseções como A.1, B.1
%%-------------------------------
\newcommand{\startappendix}{
	\appendix
	\setcounter{section}{0}                   
	\renewcommand{\thesection}{\Alph{section}} 
	\renewcommand{\thesubsection}{\thesection.\arabic{subsection}}
}
%%-------------------------------
%% Cor cinza para tabelas
%%-------------------------------

